% Part: first-order-logic
% Chapter: completeness
% Section: compactness-direct

\documentclass[../../../include/open-logic-section]{subfiles}

\begin{document}

\iftag{FOL}
      {\olfileid[id]{fol}{com}{cpd}}
      {\olfileid[id]{pl}{com}{cpd}}

\olsection{Pembuktian Langsung Teorema Kekompakan}

Kita dapat membuktikan teorema kekompakan secara langsung, tanpa
menggunakan teorema kelengkapan, dengan gagasan yang sama seperti dalam
pembuktian teorema kelengkapan. Dalam pembuktian teorema kelengkapan,
kita mulai dengan himpunan konsisten~$\Gamma$ dari !!{sentence},
memperluasnya menjadi himpunan !!{sentence} yang
konsisten\iftag{FOL}{, jenuh,}{} dan !!{complete}, yaitu~$\Gamma^*$,
lalu menunjukkan bahwa
dalam \iftag{FOL}{model suku~$\Struct{M(\Gamma^*)}$}{!!{valuation}
~$\pAssign{v(\Gamma^*)}$} yang dikonstruksi dari $\Gamma^*$, semua
!!{sentence} dalam~$\Gamma$ benar, sehingga $\Gamma$ dapat dipenuhi.

Kita dapat menggunakan metode yang sama untuk menunjukkan bahwa suatu
himpunan kalimat yang dapat dipenuhi secara berhingga dapat dipenuhi.
Kita hanya perlu membuktikan versi-versi yang bersesuaian dari
hasil-hasil yang menuju lema kebenaran, dengan mengganti ``konsisten''
dengan ``dapat dipenuhi secara berhingga.''

\begin{prop}
\ollabel{prop:fsat-ccs}
Andaikan $\Gamma$ !!{complete} dan dapat dipenuhi secara berhingga. Maka:
\begin{enumerate}
\tagitem{prvAnd}{$(!A \land !B) \in \Gamma$ jika dan hanya jika sekaligus
  $!A \in \Gamma$ dan $!B \in \Gamma$.}{}

\tagitem{prvOr}{$(!A \lor !B) \in \Gamma$ jika dan hanya jika
  $!A \in \Gamma$ atau $!B \in \Gamma$.}{}

\tagitem{prvIf}{$(!A \lif !B) \in \Gamma$ jika dan hanya jika
  $!A \notin \Gamma$ atau $!B \in \Gamma$.}{}
\end{enumerate}
\end{prop}

\tagprob{FOL}
\begin{prob}
Buktikan \olref[fol][com][cpd]{prop:fsat-ccs}. Hindari penggunaan
$\Proves$.
\end{prob}
\tagendprob

\tagprob{notFOL}
\begin{prob}
Buktikan \olref[pl][com][cpd]{prop:fsat-ccs}. Hindari penggunaan
$\Proves$.
\end{prob}
\tagendprob

\iftag{FOL}{%
\begin{lem}
\ollabel{lem:fsat-henkin} Setiap himpunan~$\Gamma$ yang dapat dipenuhi
secara berhingga dapat diperluas menjadi himpunan~$\Gamma'$ yang jenuh
dan dapat dipenuhi secara berhingga.
\end{lem}
}{}

\tagprob{FOL}
\begin{prob}
Buktikan \olref[fol][com][cpd]{lem:fsat-henkin}. (Petunjuk: langkah
pentingnya ialah menunjukkan bahwa jika $\Gamma_n$ dapat dipenuhi secara
berhingga, maka $\Gamma_n \cup \{!D_n\}$ juga demikian, tanpa menggunakan
!!{derivation} ataupun konsistensi.)
\end{prob}
\tagendprob

\iftag{FOL}{
\begin{prop}\ollabel{prop:fsat-instances}
Andaikan $\Gamma$ lengkap, dapat dipenuhi secara berhingga, dan jenuh.
\begin{tagenumerate}{prvEx,prvAll}
\tagitem{prvEx}{$\lexists[x][!A(x)] \in \Gamma$ jika dan hanya jika
  $!A(t) \in \Gamma$ untuk sekurang-kurangnya satu suku tertutup~$t$.}{}
\tagitem{prvAll}{$\lforall[x][!A(x)] \in \Gamma$ jika dan hanya jika
  $!A(t) \in \Gamma$ untuk semua suku tertutup~$t$.}{}
\end{tagenumerate}
\end{prop}
}{}

\tagprob{FOL}
\begin{prob}
Buktikan \olref[fol][com][cpd]{prop:fsat-instances}.
\end{prob}
\tagendprob

\begin{lem}
\ollabel{lem:fsat-lindenbaum} Setiap himpunan~$\Gamma$ yang dapat
dipenuhi secara berhingga dapat diperluas menjadi himpunan~$\Gamma^*$
yang !!a{complete} dan dapat dipenuhi secara berhingga.
\end{lem}

\tagprob{FOL}
\begin{prob}
Buktikan \olref[fol][com][cpd]{lem:fsat-lindenbaum}. (Petunjuk: langkah
pentingnya ialah menunjukkan bahwa jika $\Gamma_n$ dapat dipenuhi secara
berhingga, maka salah satu dari $\Gamma_n \cup \{!A_n\}$ dan $\Gamma_n
\cup \{\lnot !A_n\}$ dapat dipenuhi secara berhingga.)
\end{prob}
\tagendprob

\tagprob{notFOL}
\begin{prob}
Buktikan \olref[pl][com][cpd]{lem:fsat-lindenbaum}. (Petunjuk: langkah
pentingnya ialah menunjukkan bahwa jika $\Gamma_n$ dapat dipenuhi secara
berhingga, maka salah satu dari $\Gamma_n \cup \{!A_n\}$ dan $\Gamma_n
\cup \{\lnot !A_n\}$ dapat dipenuhi secara berhingga.)
\end{prob}
\tagendprob

\begin{thm}[Kekompakan]
\ollabel{thm:compactness-direct} $\Gamma$ dapat dipenuhi jika dan hanya
jika dapat dipenuhi secara berhingga.
\end{thm}

\begin{proof}
Jika $\Gamma$ dapat dipenuhi, maka terdapat
\iftag{FOL}{!!a{structure}~$\Struct{M}$}{!!a{valuation}~$\pAssign{v}$}
sedemikian sehingga $\iftag{FOL}{\Sat{M}}{\pSat{v}}{!A}$ untuk setiap
$!A \in \Gamma$. Tentu saja,
\iftag{FOL}{$\Struct{M}$}{$\pAssign{v}$} ini juga memenuhi setiap
himpunan bagian berhingga dari~$\Gamma$, sehingga $\Gamma$ dapat dipenuhi
secara berhingga.

Sekarang andaikan $\Gamma$ dapat dipenuhi secara berhingga.\iftag{FOL}{
  Berdasarkan \olref{lem:fsat-henkin}, terdapat himpunan $\Gamma'
  \supseteq \Gamma$ yang dapat dipenuhi secara berhingga dan jenuh.}{}
Berdasarkan \olref{lem:fsat-lindenbaum},
$\iftag{FOL}{\Gamma'}{\Gamma}$ dapat diperluas menjadi himpunan
$\Gamma^*$ yang !!a{complete} dan dapat dipenuhi secara
berhingga\iftag{FOL}{, serta $\Gamma^*$ tetap jenuh}{}. Konstruksikan
\iftag{FOL}{model suku~$\Struct{M(\Gamma^*)}$}{!!{valuation}
~$\pAssign{v(\Gamma^*)}$} seperti dalam
\olref[mod]{defn:termmodel}.\iftag{FOL}{ Perhatikan bahwa
\olref[mod]{prop:quant-termmodel} tidak bergantung pada fakta bahwa
$\Gamma^*$ konsisten (atau !!{complete} ataupun jenuh), melainkan hanya
pada fakta bahwa $\Struct{M(\Gamma^*)}$ tercakup.}{} Pembuktian Lema
Kebenaran (\olref[mod]{lem:truth}) tetap berlaku jika rujukan pada
\olref[ccs]{prop:ccs} diganti dengan rujukan pada
\olref{prop:fsat-ccs}\iftag{FOL}{, dan rujukan pada
\olref[hen]{prop:saturated-instances} diganti dengan rujukan pada
\olref{prop:fsat-instances}}.
\end{proof}

\tagprob{FOL}
\begin{prob}
Tuliskan pembuktian lengkap Lema Kebenaran
(\olref[fol][com][mod]{lem:truth}) dalam versi yang diperlukan bagi
pembuktian \olref[fol][com][cpd]{thm:compactness-direct}.
\end{prob}
\tagendprob

\tagprob{notFOL}
\begin{prob}
Tuliskan pembuktian lengkap Lema Kebenaran
(\olref[pl][com][mod]{lem:truth}) dalam versi yang diperlukan bagi
pembuktian \olref[pl][com][cpd]{thm:compactness-direct}.
\end{prob}
\tagendprob

\end{document}
